\begin{table}[htbp]
\centering
\small
\begin{threeparttable}
\caption{Standardized Effect Sizes}
\label{tab:sde}
\begin{tabular}{lcccccp{2.5cm}}
\toprule
Outcome & $\hat{\beta}$ & SE & SD($Y$) & SDE & SE(SDE) & Classification \\
\midrule
\multicolumn{7}{l}{\textit{Panel A: Pooled}} \\[3pt]
Log Wages (D11) & 0.391 & 0.029 & 0.193 & 2.024 & 0.152 & Large positive \\
Log Non-Wage Costs & -1.289 & 0.056 & 0.451 & -2.860 & 0.124 & Large negative \\
Non-Wage Share & -0.348 & 0.017 & 0.123 & -2.835 & 0.137 & Large negative \\
\midrule
\multicolumn{7}{l}{\textit{Panel B: Heterogeneous (Tradeable vs. Non-Tradeable Sectors)}} \\[3pt]
Log Non-Wage (Tradeable) & -1.164 & 0.047 & 0.414 & -2.814 & 0.113 & Large negative \\
Log Non-Wage (Non-Tradeable) & -1.337 & 0.050 & 0.489 & -2.735 & 0.103 & Large negative \\
\bottomrule
\end{tabular}
\end{threeparttable}
\vspace{0.3em}
\parbox{\textwidth}{\footnotesize
\item \textit{Notes:} \textbf{Country:} Romania and five CEE comparison countries (Bulgaria, Czechia, Hungary, Poland, Slovakia). \textbf{Research question:} Does an overnight shift of employer social security contributions to employees affect the composition of labor costs, and does statutory incidence determine economic incidence? \textbf{Policy mechanism:} Romania's GEO 79/2017 transferred nearly all employer social contributions (from 22.75\% to 2.25\%) to employees (from 16.5\% to 35\%), accompanied by a mandatory gross wage increase to offset net pay effects---the largest single-day statutory incidence reversal in modern European history. \textbf{Outcome definition:} Eurostat Labour Cost Index components: D11 (wages and salaries), D12\_D4\_MD5 (employer social contributions and non-wage costs), and their ratio (non-wage share of total compensation). All indices seasonally and calendar adjusted, base 2020$=$100. \textbf{Treatment:} Binary; Romania treated at 2018-Q1, CEE peers as controls. \textbf{Data:} Eurostat LCI (lc\_lci\_r2\_q), quarterly, 6 countries, aggregate business sector B--S, 2016-Q1 to 2019-Q4 (16 quarters). \textbf{Method:} Two-way fixed effects DiD (country + quarter FE), standard errors clustered at country level. Permutation inference yields $p < 0.05$ for both wage and non-wage outcomes. \textbf{Sample:} Six CEE EU member states with complete LCI data 2016--2019; stopped at 2019-Q4 to avoid COVID contamination. SDE $= \hat{\beta} / \text{SD}(Y)$ where SD($Y$) is the pre-treatment standard deviation. Classification refers to magnitude, not statistical significance: Large ($|$SDE$| > 0.15$), Moderate ($0.05$--$0.15$), Small ($0.005$--$0.05$), Null ($< 0.005$).
}
\end{table}
